\chapter*{Conclusion: Sharing Time Without Sharing Every Frame}
\addcontentsline{toc}{chapter}{Conclusion: Sharing Time Without Sharing Every Frame}
\markboth{Conclusion}{Conclusion}

Return to the two people in adjacent seats. They came in together, sat in the dark, and looked at the same screen for an hour and a half. The film began for both at once and ended for both at once. They heard the same music, the same doors slamming, the same laughter from the far side of the room. They followed the same characters through the same events, and at several moments in the film, though neither knew exactly when, everything each of them had been shown agreed on everything that would matter afterward. One of them saw a thriller and the other a comedy. In the last minutes their films resolved in opposite directions, and on the way out the first thing they asked each other was which one they had seen.

The book has tried to say what makes that evening possible and what makes it worth having.

The apparatus is modest. A projector already shows more flashes than one film needs, and glasses that open on a chosen period and phase can divide those flashes among several films. The division is costly, in light and in flicker, and the costs rise with every stream, so the families that can be shown are small. But the costs fall whenever the films coincide, and the room without glasses sees their average, which can itself be designed.

The form is demanding. A family of films that can share one screening must be written together, against a shared spine, with scenes that meet on a common clock and moments at which every realization agrees on what has happened. Whether a viewer can move from one realization to another is not a matter of taste but of continuity: whether their history can be carried into the other's future without leaving them lost or deceived. Sound, which cannot be divided the way light can, is best composed once for all the realizations, so that the room stays acoustically whole while the pictures diverge.

The stakes are real. A screen that shows different people different things can widen what a film can do: offer comedy and drama, beginner and expert, one ending and another, access and translation, to people who remain together. It can also narrow what a public can know, if the variation is hidden, assigned without consent, sold by tier, locked behind keys in devices the viewers cannot inspect, or used to deliver different messages under one title. The difference between these outcomes is not in the optics. It is in whether the variation is declared, chosen, and made by people with standing to make it.

The central claim can now be stated with the precision the book has earned. Collective viewing does not require that everyone see the same thing. It requires that different experiences stay bounded in time, return repeatedly to agreement, and share at least one channel that everyone receives alike. Where those conditions hold, an audience is neither a mass receiving one message nor a scatter of individuals receiving unrelated ones. It is a plurality that shares time without sharing every frame.

That arrangement has been present in the idea from its first form: two films about a society whose rulers controlled what everyone saw, made to be watched at once, joined by a score written for both, and separated only by the timing of a shutter each viewer set for themselves. The story was about the danger of one authority deciding what a public sees. The way it was to be shown was an answer to that danger: not one picture for everyone, and not a private picture for each, but several pictures, openly offered, freely chosen, and watched together.
